
\section{Finite convex representations}\label{sec:atoms}

We next describe the finite-dimensional dual balls. Exact finite
representations will let us perform the cancellation calculation with
finite probability spaces, postponing the only limit until a coordinate
exhaustion at the end.

For a segment $S$, write $d_S=\sum_{v\in S}e_v^*$. By
\eqref{eq:segment-norm},
\begin{equation}\label{eq:atom-coefficients}
 \left\|\sum_i a_i d_{S_i}\right\|
 \le\left(\sum_i a_i^2\right)^{1/2}
 \qquad(S_i\text{ pairwise disjoint}).
\end{equation}
Indeed, evaluate on a vector $u$ with $\rho(u)\le1$ and apply
Cauchy--Schwarz to the coefficient vector and the segment sums. Call
such a represented functional an \emph{atom} when $\sum_i a_i^2\le1$.
The representation, including the segment family, is part of the data.

\begin{lemma}\label{lem:finite-hull}
For a finite ancestral set $K$, the unit ball of $P_KX$ is the convex
hull of the atoms whose segments lie in $K$. Each member is a convex
combination of at most $|K|+1$ such atoms. If $0\ne x=P_Hx$ for a
finite ancestral $H$, there is $p=P_Hp\in J$ with $\rho(p)=1$ and
$x(p)=\|x\|$.
\end{lemma}

\begin{proof}
There are finitely many disjoint segment families in $K$. For each,
its atoms are the continuous image of a Euclidean unit ball, so the
union is compact. Its convex hull is compact in the finite-dimensional
space $P_KX$. For a test $p$ supported in $K$, optimizing the atom
coefficients shows that the support function of this convex hull is
$\rho(p)$. In computing $\rho(p)$, all segment families can be
restricted to $K$. Finite-dimensional separation therefore identifies
the hull with the dual unit ball on that coordinate space. By
contractivity of $P_K$, this is precisely the inherited unit ball of
$P_KX$.

To obtain the stated bound on the number of atoms, consider a convex
representation with more than $|K|+1$ positive weights. Its atoms are
affinely dependent, so there are real numbers $\theta_i$, not all zero,
with $\sum_i\theta_i=0$ and $\sum_i\theta_i A_i=0$. Vary the weights
along this direction until the first weight reaches zero; all weights
remain nonnegative, and both their sum and their represented vector
are unchanged. Repeating leaves at most $|K|+1$ terms.

Likewise, for $x=P_Hx$ the dual pairing may be restricted to $P_HJ$
without altering its norm. That space is finite dimensional, so its
compact unit ball contains a norming test, of norm one when $x\ne0$.
\end{proof}

\begin{corollary}[The completed atomic ball]\label{cor:closed-atomic-ball}
Let $\mathcal A$ be the set of all represented atoms in $X$. Then
\[
 B_X=\overline{\operatorname{conv}\mathcal A}^{\,\|\cdot\|}.
\]
Every segment functional $d_S$ has norm one.
\end{corollary}

\begin{proof}
The coefficient estimate puts every atom in $B_X$. Conversely, for
$x\in B_X$, choose an increasing finite ancestral exhaustion $(K_n)$.
Contractivity gives $P_{K_n}x\in B_X$, and
Lemma~\ref{lem:finite-hull} puts each $P_{K_n}x$ in
$\operatorname{conv}\mathcal A$. The convergence
$P_{K_n}x\to x$ proves the reverse inclusion. For a nonempty segment
$S$, its functional has norm at most one by the coefficient estimate,
and evaluation on $e_v$, for any $v\in S$, gives equality: the singleton
vector $e_v$ has $\rho(e_v)=1$.

There is also a direct separation proof of the completed hull identity.
Optimizing the atom coefficients gives
$\sup_{A\in\mathcal A}A(p)=\rho(p)$ for every $p\in J$.
Since $X^*=J$, separation of any point of $B_X$ from the norm-closed
convex hull would give $x(p)>\rho(p)$, contradicting $\|x\|\le1$.
This second proof uses reflexivity; the finite-projection proof only
uses the coordinate completion and its contractive exhaustion.
\end{proof}

\section{Cancellation at the boundary}\label{sec:cancellation}

We can now compare the segment coefficients of two endpoint laws. The
first step depends only on the head test, not on the original atom
coefficients. It explains why opposite crossing sums cannot be large
when both segment tests are nearly maximal.

For a finite ancestral $H$, a segment \emph{crosses at $g\in G_H$} if
it meets $H$ and its first vertex outside $H$ is $g$. For a disjoint
segment family $\mathcal S$ and a test $p=P_Hp$ with $\rho(p)\le1$,
put
\[
 L_{\mathcal S}=1-\sum_{S\in\mathcal S}p(S)^2,\qquad
 b_{\mathcal S}(g)=
 \begin{cases}
 p(S),&S\in\mathcal S\text{ crosses at }g,\\
 0,&\text{no segment of }\mathcal S\text{ crosses at }g.
 \end{cases}
\]
At most one segment of a disjoint family crosses at each gate, so this
definition is unambiguous. Also $L_{\mathcal S}\ge0$.

\begin{lemma}\label{lem:opposite-crossing}
With the preceding notation, two disjoint segment families
$\mathcal S,\mathcal S'$ satisfy
\begin{equation}\label{eq:opposite-crossing}
 2\sum_{g:\,b_{\mathcal S}(g)b_{\mathcal S'}(g)<0}
   |b_{\mathcal S}(g)b_{\mathcal S'}(g)|
 \le L_{\mathcal S}+L_{\mathcal S'}.
\end{equation}
\end{lemma}

\begin{proof}
At a gate with opposite nonzero sums $b,b'$, the two heads end at the
same parent $v$. They are therefore nested and cannot coincide.
Suppose the unprimed head $[a,v]$ is longer than the primed head
$[a',v]$. Split the unprimed segment immediately before $a'$, retaining
its own endpoint in the second piece. Since $p$ vanishes outside $H$,
the two new sums are $b-b'$ and $b'$. The gain in its squared test is
\[
 (b-b')^2+(b')^2-b^2=2(b')^2-2bb'\ge2|bb'|.
\]
Figure~\ref{fig:opposite-crossing} distinguishes the two original
segments from the two pieces created by the split.
\begin{figure}[tb]
\centering
\begin{tikzpicture}[x=1cm,y=0.75cm,font=\small,
  head/.style={line width=1.4pt,blue!65!black},
  second/.style={line width=1.4pt,red!65!black},
  tail/.style={dashed,line width=1pt},
  vertex/.style={circle,fill=black,inner sep=1.5pt}]
  \node at (1.4,4.6) {Two crossing heads};
  \draw[gray!70,densely dotted] (-0.6,0.4)--(3.1,0.4);
  \node[anchor=west] at (2.6,0.85) {\(H\)};
  \node[anchor=west] at (2.6,-0.05) {\(H^c\)};
  \draw[head] (0,3.8)--(0,0.8);
  \draw[head,tail] (0,0.8)--(0,-0.7);
  \draw[second] (1.1,2.2)--(1.1,0.8);
  \draw[second,tail] (1.1,0.8)--(1.1,-0.7);
  \node[vertex,label=left:\(a\)] at (0,3.8) {};
  \node[vertex,label=left:\(a'\)] at (0,2.2) {};
  \node[vertex] at (1.1,2.2) {};
  \node[vertex,label=left:\(v\)] at (0,0.8) {};
  \node[vertex] at (1.1,0.8) {};
  \node[vertex,label=left:\(g\)] at (0,0) {};
  \node[vertex] at (1.1,0) {};
  \node[anchor=west,text=blue!65!black] at (0.15,3.1) {\(b\)};
  \node[anchor=west,text=red!65!black] at (1.25,1.5) {\(b'\)};
  \node at (1.1,-1.3) {\(bb'<0\)};
  \draw[-{Stealth[length=5pt]},gray!80] (3.5,2.1)--(4.7,2.1);
  \node[align=center] at (4.1,3.05) {Split before\\\(a'\)};
  \begin{scope}[xshift=6cm]
    \node at (1.3,4.6) {Splitting the longer segment};
    \draw[gray!70,densely dotted] (-0.6,0.4)--(3.2,0.4);
    \draw[head] (0,3.8)--(0,2.7);
    \draw[second] (0,2.2)--(0,0.8);
    \draw[second,tail] (0,0.8)--(0,-0.7);
    \node[vertex,label=left:\(a\)] at (0,3.8) {};
    \node[vertex] at (0,2.7) {};
    \node[vertex,label=left:\(a'\)] at (0,2.2) {};
    \node[vertex,label=left:\(v\)] at (0,0.8) {};
    \node[vertex,label=left:\(g\)] at (0,0) {};
    \node[anchor=west,text=blue!65!black] at (0.25,3.3)
      {sum \(b-b'\)};
    \node[anchor=west,text=red!65!black] at (0.25,1.5)
      {sum \(b'\)};
    \node at (1.25,-1.3) {Gain \(\ge2|bb'|\)};
  \end{scope}
\end{tikzpicture}
\caption{Opposite $p$-sums at a gate $g$, with $p$ supported in $H$.
The two heads on the left belong to different segment families;
their traces are displaced to show their nesting on the same ancestral
chain. Their tails are schematic and need not coincide.
On the right, the longer segment is split before $a'$; the lower
piece retains that segment's original tail and has $p$-sum $b'$.
The squared segment test increases by at least $2|bb'|$.}
\label{fig:opposite-crossing}
\end{figure}


Do this at every opposite-sign gate, in the family with the longer
head. Within each family each original segment is cut at most once,
because it crosses only one gate. Disjointness is preserved, so the
total gain in each family is at most its deficit $L$. Summing proves
\eqref{eq:opposite-crossing}.
\end{proof}

